\capitulo{6}{Trabajos relacionados}

No se han encontrado trabajos relacionados que integren de forma conjunta las
funcionalidades propuestas en este \acrshort{tfg}, combinando técnicas de
modelado de temas con procesos de ingesta y análisis de información textual
obtenida desde repositorios de GitHub.

Sin embargo, sí existen diversos \acrlong{tfg} relacionados parcialmente con
algunas de las áreas abordadas en este proyecto, especialmente en el ámbito de
la extracción de información desde GitHub, la minería de texto y el
procesamiento del lenguaje natural.

En relación con la extracción y análisis de información procedente de
repositorios GitHub, se han identificado varios trabajos desarrollados en
cursos anteriores.

El primero de ellos, presentado en junio de 2019, lleva por título
\textit{Autoevaluación del proceso de desarrollo de Trabajo Fin de Grado en
      GitHub}, desarrollado por Roberto Luquero y disponible en el repositorio
\url{https://github.com/rlp0019/Activiti-Api}. Este proyecto extrae métricas
cuantitativas obtenidas tanto del sistema de control de versiones Git como del
sistema de gestión de incidencias (\textit{issues}), utilizando posteriormente
dichos valores para compararlos con umbrales predefinidos con fines de
evaluación académica.

Posteriormente, en junio de 2022, Joaquín García Molina presentó el trabajo
\textit{Comparador de métricas de evolución en repositorios software},
disponible en
\url{https://github.com/Joaquin-GM/GII_O_MA_19.07-Comparador-de-metricas-de-evolucion-en-repositorios-Software}.
Este proyecto sigue una aproximación similar, centrada en la obtención y
comparación de métricas cuantitativas asociadas a la evolución de repositorios
software mediante información procedente de Git y del sistema de incidencias.

En junio de 2023 se desarrolló el TFG \textit{Evaluación de simulaciones de
      proyectos en GitHub orientadas a docencia}, realizado por Licinio Fernández
Maurelo y disponible en \url{https://github.com/lfx1001/}. A diferencia de los
trabajos anteriores, este proyecto se centra en el análisis completo del ciclo
de vida asociado a la resolución de incidencias, comparando distintos flujos de
trabajo para determinar el grado de similitud entre ellos.

Los tres trabajos anteriores comparten un enfoque orientado al ámbito docente,
proporcionando herramientas de apoyo para que los estudiantes puedan evaluar y
comparar sus procesos de desarrollo realizados en GitHub frente a proyectos
considerados de referencia o de alta calidad.

En el contexto de la minería de texto, en enero de 2022 se presentó el proyecto
\textit{Plataforma de experimentación de text mining sobre repositorios de
      código abierto GitHub}, desarrollado por Pablo Fernández Bravo y disponible en
\url{https://github.com/2021-22-TFG-ene/github-text-mining-tfg}. Este trabajo
emplea la información contenida en las incidencias de los repositorios para
realizar tareas de clasificación automática de texto, análisis de sentimientos
y generación de resúmenes.

Finalmente, dentro del ámbito del resumen automático de textos mediante
técnicas de procesamiento del lenguaje natural, en febrero de 2021 se presentó
el TFG \textit{JIZT. Servicio de Resumen de Textos con AI en la Nube},
desarrollado por Diego de Miguel Lozano y disponible en
\url{https://github.com/dmlls/jizt-tfg}. Este proyecto se centra en la
generación automática de resúmenes textuales mediante técnicas de inteligencia
artificial aplicadas al procesamiento del lenguaje natural.

\begin{table}[H]
      \centering
      \caption{Resumen de trabajos relacionados}
      \label{tab:trabajos_relacionados}
      \begin{tabular}{|c|p{4.5cm}|p{3cm}|p{3cm}|}
            \hline
            \textbf{Año} & \textbf{Trabajo}                                                                                  & \textbf{Área principal}                       & \textbf{Relación con el TFG}                              \\
            \hline 2019  & \textit{Autoevaluación del proceso de desarrollo de Trabajo Fin de Grado en GitHub}               & Métricas sobre Git e incidencias              & Extracción y análisis cuantitativo de repositorios GitHub \\
            \hline 2021  & \textit{JIZT. Servicio de Resumen de Textos con AI en la Nube}                                    & Resumen automático de texto                   & Procesamiento del lenguaje natural                        \\
            \hline 2022  & \textit{Comparador de métricas de evolución en repositorios software}                             & Métricas de evolución software                & Comparación de métricas de repositorios                   \\
            \hline 2022  & \textit{Plataforma de experimentación de text mining sobre repositorios de código abierto GitHub} & Minería de texto                              & Clasificación, análisis de sentimientos y resumen         \\
            \hline 2023  & \textit{Evaluación de simulaciones de proyectos en GitHub orientadas a docencia}                  & Simulación y comparación de flujos de trabajo & Análisis del ciclo de vida de incidencias                 \\
            \hline
      \end{tabular}
\end{table}
